\chapter{1946:James B. Sumner and John H. Northrop and Wendell M. Stanley}

\label{chap:chemistry1946}

The Nobel Prize in Chemistry for the year 1946 was awarded jointly to James B. Sumner and John H. Northrop and Wendell M. Stanley.

\textbf{Motivation:}

for his discovery that enzymes can be crystallized (one half, James B. Sumner) and for their preparation of enzymes and virus proteins in a pure form (the other half, jointly, John H. Northrop and Wendell M. Stanley)

\section{James B. Sumner}

\subsection*{Early Life and Education}

James B. Sumner was born on 1887-11-19 in Canton, Massachusetts, USA.

\subsection*{Career and Major Research}

Sumner studied at Harvard University, where he earned his doctorate in 1914. He then joined Cornell University, where he spent his career as a professor of biochemistry. In 1926 he isolated and crystallized the enzyme urease from jack beans, proving that a purified enzyme is a protein. The result was disputed for years by chemists who believed enzymes were non-protein substances.

\subsection*{Nobel Prize-winning Work}

The work for which James B. Sumner was awarded the Nobel Prize in Chemistry in 1946 was described by the Nobel Committee as: "for his discovery that enzymes can be crystallized".

Crystallization made it possible to obtain an enzyme as a pure chemical substance and to show unambiguously that the enzyme itself was protein in nature.

\subsection*{Significance and Impact}

Sumner's result settled the question of the chemical nature of enzymes and opened the way to their purification and detailed study. It established enzyme chemistry as a field in which individual catalysts could be isolated, characterized, and compared.

\subsection*{Later Career and Honors}

Sumner continued his research at Cornell University until his retirement. He died in 1955.

\subsection*{Legacy}

The crystallization of urease is regarded as the demonstration that enzymes are proteins, a cornerstone of modern biochemistry and molecular biology.

\subsection*{Selected Publications}

\begin{itemize}

\item \emph{The Isolation and Crystallization of the Enzyme Urease}, Journal of Biological Chemistry, 1926.

\end{itemize}

\clearpage

\section{John H. Northrop}

\subsection*{Early Life and Education}

John H. Northrop was born on 1891-07-05 in Yonkers, New York, USA.

\subsection*{Career and Major Research}

Northrop studied chemistry at Columbia University and joined the Rockefeller Institute for Medical Research, where he worked for most of his career. In 1929 he crystallized the enzyme pepsin, and later crystallized the digestive enzymes trypsin and chymotrypsin. These results independently confirmed Sumner's conclusion that enzymes are proteins.

\subsection*{Nobel Prize-winning Work}

The work for which John H. Northrop was awarded the Nobel Prize in Chemistry in 1946 was described by the Nobel Committee as: "for their preparation of enzymes and virus proteins in a pure form".

Northrop's crystallization of pepsin and the other protein-digesting enzymes established beyond doubt that enzymes are proteins and could be purified to crystalline form.

\subsection*{Significance and Impact}

Northrop's work provided the decisive confirmation, from an independent laboratory, that enzymes are proteins. It also made crystalline enzyme preparations available for structural and kinetic study.

\subsection*{Later Career and Honors}

Northrop remained at the Rockefeller Institute for Medical Research, where he directed work in enzyme chemistry. He died in 1987.

\subsection*{Legacy}

The crystalline enzymes prepared by Northrop were among the first proteins to be studied as pure chemical substances, helping to found protein chemistry.

\subsection*{Selected Publications}

\begin{itemize}

\item \emph{Crystalline Pepsin}, Science, 1929.

\end{itemize}

\clearpage

\section{Wendell M. Stanley}

\subsection*{Early Life and Education}

Wendell M. Stanley was born on 1904-08-16 in Ridgeville, Indiana, USA.

\subsection*{Career and Major Research}

Stanley studied at Earlham College and the University of Illinois, where he received his doctorate in 1929. He worked at the Rockefeller Institute for Medical Research in Princeton, New Jersey. In 1935 he isolated tobacco mosaic virus in crystalline form and showed that the virus particle consisted largely of protein.

\subsection*{Nobel Prize-winning Work}

The work for which Wendell M. Stanley was awarded the Nobel Prize in Chemistry in 1946 was described by the Nobel Committee as: "for their preparation of enzymes and virus proteins in a pure form".

Stanley applied crystallization to a virus rather than an enzyme, obtaining tobacco mosaic virus as a crystalline material that retained its infectivity. The work showed that a virus could behave as a well-defined chemical substance.

\subsection*{Significance and Impact}

Stanley's crystallization of tobacco mosaic virus demonstrated that viruses, too, could be studied as pure chemical entities. It helped establish the molecular basis of virology and the protein chemistry of viruses.

\subsection*{Later Career and Honors}

Stanley later served as professor of biochemistry at the University of California, Berkeley, where he directed a laboratory of virology. He died in 1971.

\subsection*{Legacy}

Stanley's work joined the study of viruses to chemistry, and his crystalline tobacco mosaic virus became a model system in the early development of molecular biology.

\subsection*{Selected Publications}

\begin{itemize}

\item \emph{Isolation of a Crystalline Protein Possessing the Properties of Tobacco-Mosaic Virus}, Science, 1935.

\end{itemize}

\clearpage

\section*{Chapter Summary}

The 1946 prize recognized the crystallization of enzymes, which proved that enzymes are proteins. James B. Sumner crystallized urease, John H. Northrop crystallized pepsin and other digestive enzymes, and Wendell M. Stanley extended the approach to tobacco mosaic virus, showing that a virus could be obtained as crystalline material.
